\documentclass[11pt]{article}
\usepackage[letterpaper,margin=0.65in]{geometry}
\usepackage[T1]{fontenc}
\usepackage[default]{sourcesanspro}
\usepackage{amsmath,amssymb,enumitem,fancyhdr,xcolor}
\setlength{\emergencystretch}{1em}
\setlength{\parindent}{0pt}\setlength{\parskip}{6pt}
\definecolor{ink}{HTML}{183642}
\pagestyle{fancy}\fancyhf{}\fancyhead[L]{\small\textcolor{ink}{Week 9 / Bouncing paths / Adult return-visit guide}}
\fancyfoot[L]{\small Bellingham Math Circle / Week 9 / RV9-FAC-v1}\fancyfoot[R]{\small\thepage}
\renewcommand{\headrulewidth}{0.4pt}\renewcommand{\footrulewidth}{0pt}
\setlist[itemize]{leftmargin=1.2em,itemsep=2pt,topsep=2pt}
\newcommand{\parttitle}[1]{{\large\bfseries\color{ink}#1}\par}
\newcommand{\labelled}[2]{\textbf{#1} #2\par}
\begin{document}
\parttitle{Mathematical destination}
\labelled{Ordered two-bounce routes.}{In an ideal rectangle, fix distinct strictly interior S and F whose x and y coordinates both differ. This excludes normal-incidence retracing and an earlier target passage in these examples. Seek exactly two specular wall contacts, excluding corners and extra contacts. A wall word has at most one route: reflect F through its last wall, then its first, and check the straight candidate's actual wall order after folding. Same-wall words are impossible. Opposite-wall orders LR, RL, BT, TB always work. For each adjacent pair of walls, exactly one order works unless the candidate hits their corner, when neither works. Perpendicular reflections commute, but their routes' orders do not. On the displayed 4-by-4 examples there are eight and six legal words.}
\labelled{Interior starts can loop.}{On a 6-wide,4-high table with a northeast 45-degree start at an interior lattice dot $(x,y)$, a corner is reached iff $x-y$ is even. Of the 15 dots, eight reach corners and seven never do. The latter return to the same dot travelling northeast after 24 unit diagonal steps; returning to a dot with a different direction does not finish the experiment.}
\labelled{Crossing geometry.}{For an ordinary 45-degree corner start, reduce dimensions by $g=\gcd(a,b)$ to coprime $u=a/g$, $v=b/g$. There are $(u-1)(v-1)/2$ distinct interior transverse crossing points before the first terminal corner. No segment is retraced. Enlarging both dimensions by the same factor preserves the count. Count a crossing location once, excluding walls and terminal corners.}
\labelled{Discoveries versus guarantees.}{Children draw exact paths and compare wall words, starts and crossing points. A drawing establishes that example. The mirror-copy arguments explain all rectangles under the stated assumptions; a finite table alone does not prove a general formula. These investigations add ordered inverse targeting, state and intersection geometry to base Week 9's endpoint/bounce/lcm work.}
\labelled{Band map.}{Grades 2--3 can test two-bounce routes from an explicit reflection visual and mark crossing points. Interior-start pages are marked 4--5; a ready third grader can try with an adult checking the visible direction arrow. Grades 4--5 can explain unfolding, parity and scaling. A K--1 floor walk with an adult is possible, but no independent younger version is claimed. Ruler use, grid tracing and following a slope are the concrete prerequisites; lcm and formula manipulation belong to later explanations.}
\labelled{Status.}{Prepared and unpiloted; exact material and floor-walk procedures have not been physically rehearsed. Give one investigation time to develop and keep unfinished paths for return visits.}
\newpage
\parttitle{Prepare and launch}
\labelled{Per pair.}{Ruler, two or three coloured pencils, scraps of square-grid paper, a small token and four direction-arrow cards NE,SE,SW,NW if walking a path. Use equally scaled grid axes. Two-bounce and crossing grids use 12.5-mm pencil squares. The main interior-start board uses 22-mm units, and its workspace boards use 18-mm units. A tracking token at most 12 mm across is a sensible choice; physical fit has not been rehearsed. The arrow card can remain beside the grid. Enlarge or tape a floor grid for movement. For two-bounce work provide tracing paper and spare paper; string can compare broken route lengths in a continuation. For Problem 1 children choose the initial direction. Problems 2--3 use the specified northeast 45-degree launch.}
\labelled{Exact rules.}{At a vertical wall reverse the horizontal direction only; at a horizontal wall reverse the vertical direction only. Corner bounces are not allowed in the two-bounce task. In the ordinary diagonal-path tasks, stop at a corner. For interior-loop tasks stop on the original dot only when travelling in the original direction. Do not treat a crossing of the drawn path as a wall, turn or stop. A bounce is a wall contact before the terminal corner.}
\labelled{Launch together, 3--5 minutes.}{Allow children to move a token on a spare grid, then act out a legal reflection and let a partner check it. Before the two-bounce pages, use the non-task visual S=(1,1), R-wall contact (4,2), F=(1,3). Its reflected outgoing vector reverses horizontal direction and retains vertical direction; both angles to the wall are equal. The completed route gives wall word R. The task then defines LR as left followed by right. Let a child check the equal reflection using a fold or reflected copy. The printed diagonal convention shows (2,1) NE to (3,2) NW to (2,3) SW on a non-task 3-by-3 board, each legal turn reversing only the wall-normal direction. Before an interior-start page, place the arrow beside the token so its state includes direction. The token is a tracker on a pencil grid; no counter-fit claim is made. These are meaningful intermediate actions, not merely an unexplained finished path. Use the non-task printed example before the first unfamiliar convention.}
\labelled{A manageable route.}{Student pages 1--3 belong to Problem 1, pages 4--5 to the older loop entry, and pages 6--7 to crossings. For a grades 2--3 route, choose the two-bounce or crossing pages; give their matching launch before handing them out. Use 20--30 minutes for several wall words or crossing-point tables. A return visit can introduce interior dots and spend another 20--30 minutes tracing disputed starts. Trace a path with direction arrows or save its ordered wall contacts; count its diagonal unit steps afterward. At 45 degrees, each straight segment contributes the absolute horizontal change in grid units. Keeping this recoverable record is more helpful than simultaneous bounce, time and crossing tallies. A pair can share the 15 starts across visits; neither child needs seven concurrent loop counts. The upper table can rotate tracer, predictor and checker.}
\labelled{Adult support.}{Read labels or draw a ruler line at the child's chosen direction. If an adult must choose every turn, return to a smaller concrete walk. Let children retain the choice of start, direction and conjecture. Optional formula work need not interrupt a child absorbed in drawings.}
\labelled{Lineage.}{Lowell Fall 2025 \emph{Handout 10}, Problems 10.2--10.7, already used 45-degree tables, endpoints and scaling. Current Week 9 F09-v4 develops the general endpoint/bounce rule through unfolded copies. The archived grades 6--7 extra also used rational slopes and one midpoint periodic orbit. Week 21 already contains shortest one-contact mirror routes, so the new investigation instead prescribes two wall contacts and tests their order. Its return-visit companion also optimizes two contacts on parallel walls. Both use repeated reflection; Week 9 instead enumerates rectangle wall words and rejects wrong orders and corner hits. See the Week 21 bonus companion for the length-optimization direction. Complete interior-dot classification is another continuation; neither reflections nor loops are claimed as new themes. The constructions and exact checks here are original. \emph{Math Circle by the Bay}, preface printed viii--x/PDF 9--11, supports manipulatives, explanations and themes at different depths; our staffing, routes and times are adaptations. No classroom results are available.}
\newpage
\parttitle{Problem 1 / Two wall contacts}
\labelled{Find a candidate, then check it.}{Let L,R,B,T mean left,right,bottom,top walls. For word XY, reflect the target across wall Y, then across wall X. The line from S to that image is the only candidate: every legal XY route unfolds to it. Fold the segments back, checking that it reaches X then Y, avoids corners and reaches F without further wall contacts. A line to an image is not alone a legal route. Length is preserved under each reflection.}
\labelled{Checked answer lists.}{On the 4-by-4 table with S=$(1,1)$ and F=$(2,3)$ (Finish A), the legal words are LR, LT, RL, RT, BL, BR, BT, TB. With F=$(3,3)$ (Finish B) they are LR, LT, RL, BR, BT, TB. Thus all sixteen words, including LL,RR,BB,TT, have been checked. For the second target the LB/BL image is $(-3,-3)$ and the ray hits $(0,0)$; the RT/TR image is $(5,5)$ and the ray hits $(4,4)$. These are excluded corner contacts.}
\labelled{Why equal images need an order test.}{Reflecting across perpendicular walls in either order produces the same image. The ray crosses the supporting wall lines at distinct times, selecting one order. The opposite word has the wrong first wall. If both crossings are simultaneous, it hits a corner and neither word is legal. In the first example LT works and TL does not; BL works and LB does not.}
\labelled{Why opposite orders work.}{Reflecting across opposite walls translates the target by twice the width or height. The ray crosses those two successive wall copies before ending in the target copy. Its other coordinate stays between the corresponding coordinates of S and F, so it meets no transverse wall. The two orders give different translations and different legal routes. A repeated wall is impossible: after its reflection the ray travels away from that wall and must hit another wall before it can return.}
\labelled{Exact checks for the first target.}{Contact pairs are LR: $(0,9/7)$ then $(4,17/7)$; RL: $(4,5/3)$ then $(0,23/9)$; BT: $(7/6,0)$ then $(11/6,4)$; TB: $(13/10,4)$ then $(17/10,0)$. Adjacent words are LT: $(0,7/3)$ then $(5/4,4)$; BL: $(1/4,0)$ then $(0,1/3)$; RT: $(4,17/5)$ then $(13/4,4)$; BR: $(9/4,0)$ then $(4,7/5)$. Fractions are adult checks, not entry requirements.}
\labelled{Hints and continuation.}{Begin with one chosen word and test a proposed first contact using a ruler. If needed, give a mirror-copy target and let the child choose the straight line. Keep a single wall word and path as the record. Once legal routes are found, compare lengths using string: their squared lengths are LR 53, RL 85, BT 37, TB 101, LT 25, BL 25, RT 41, BR 41, so LT and BL tie for shortest in that family. Another return visit can move the target to make a candidate hit a corner. Real rolling and friction are outside this ideal model.}
\newpage
\parttitle{Problem 2 / Which interior dots finish?}
\labelled{Exact classification.}{On the 6-by-4 table, the corner-finishing dots are $(1,1)$, $(1,3)$, $(2,2)$, $(3,1)$, $(3,3)$, $(4,2)$, $(5,1)$, $(5,3)$. The looping dots are $(1,2)$, $(2,1)$, $(2,3)$, $(3,2)$, $(4,1)$, $(4,3)$, $(5,2)$. On the printed letter map the corner starts are A,C,E,G,I,K,M,O; the looping starts are B,D,F,H,J,L,N. Every looping dot returns northeast after 24 diagonal unit steps. The letters run left-to-right by rows: A--E at height 3, F--J at height 2, K--O at height 1. A path that passes its dot in another direction continues.}
\labelled{Why parity decides.}{In unfolded copies a northeast ray from $(x,y)$ is $(x+t,y+t)$. To hit a corner it must satisfy $t\equiv-x\pmod6$ and $t\equiv-y\pmod4$. These conditions have a solution exactly when $x-y$ is divisible by $\gcd(6,4)=2$. Necessity follows because the two left sides have the same parity. Sufficiency can be checked by comparing the even or odd residues over one 12-step interval, or by the elementary common-residue argument. The time to a finishing corner is the least positive solution, not necessarily 24.}
\labelled{Why loop time is 24.}{After reflection, a coordinate repeats with the same direction every $2a$ units. Here the horizontal period is 12 and the vertical period 8. Both coordinate directions repeat first after $\operatorname{lcm}(12,8)=24$. For an interior coordinate, a return with the same direction requires this full period, not its mirror half. A non-corner path continues until that state returns.}
\labelled{An elementary explanation route.}{Draw mirror copies of the table and mark every reflected copy of the start dot. Seek one whose outgoing arrow matches the original. Children can first compare the actual paths of two contrasting dots; they need not solve simultaneous congruences. A direction card externalizes the state. The classification is an exhaustively checked fact; a proposed general rule for other rectangles remains a conjecture until explained.}
\labelled{Hint and continuation.}{If the start-dot rule is unclear, ask whether a traveller returning home facing the other way has returned to the same state. At H=(3,2), the traveller is back at its dot after 12 steps facing SE, and first returns NE after 24; this optional checked hint illustrates why position alone is insufficient. Later try a different rectangle: the corner condition becomes $x-y$ divisible by $\gcd(a,b)$ for integer northeast starts, and a non-corner full-state return is $2\operatorname{lcm}(a,b)$. A grid dot can lead to a corner or a loop; these statements do not cover arbitrary irrational starting coordinates. Save the start and arrow with the path.}
\newpage
\parttitle{Problem 3 / Crossing locations}
\labelled{Checked counts.}{The ordinary 45-degree paths on 2-by-3, 3-by-4 and 4-by-5 tables have respectively 1,3 and 6 interior crossing points. A 4-by-6 table has one: it is the doubled 2-by-3 path. For the ten-crossing invention, 5-by-6 works, as does its doubled 10-by-12 version. With positive integer dimensions at most 12, all solutions are 3-by-11, 11-by-3, 5-by-6, 6-by-5, 10-by-12 and 12-by-10. Counts are locations, not the two visits to each location. A wall bounce and a terminal corner are excluded.}
\labelled{Explicit coordinates.}{For 2-by-3 the point is $(1,1)$. For 3-by-4 they are $(1,1)$, $(1,3)$, $(2,2)$. For 4-by-5 they are $(1,1)$, $(1,3)$, $(2,2)$, $(2,4)$, $(3,1)$, $(3,3)$. In 4-by-6 the single point is $(2,2)$. These coordinates give an independent way to inspect a drawing; do not print them as a child procedure.}
\labelled{Why the reduced count works.}{First divide the whole drawing by $g=\gcd(a,b)$; this changes scale but neither reflections nor crossing count. In the coprime $u$-by-$v$ rectangle, opposite-slope segments meet at interior lattice points $(i,j)$ with $i,j$ of the same parity. Reflected slope $+1$ or $-1$ lines have even integer intercepts, so no other crossing coordinates are possible. For a visit to $(i,j)$ the unfolded time obeys $t\equiv\pm i\pmod{2u}$ and $t\equiv\pm j\pmod{2v}$. Since $\gcd(2u,2v)=2$, matching parity gives one time modulo $2uv$ for each of the four sign pairs. Times $t$ and $2uv-t$ are paired across the full period, leaving exactly two visits before the first corner at $uv$. One visit has slope $+1$ and the other slope $-1$: the same-sign pair gives the former and the opposite-sign pair the latter. Thus no interior point is revisited along the same slope, so no segment is retraced.}
\labelled{Counting the points.}{There are $(u-1)(v-1)$ interior reduced-grid points. At least one of $u-1,v-1$ is even, since coprime $u,v$ cannot both be even. In that coordinate, equal numbers are odd and even; exactly half of all interior pairs have matching parity. Thus the count is $(u-1)(v-1)/2$. The argument also gives zero when either reduced dimension is one.}
\labelled{Discovery route and hint.}{Trace in one colour, circle crossing locations afterward and compare with a scaled table. If children count each visit twice, place one counter at each physical crossing location. If the formula is not yet interesting, continue designing a table with many crossings; the concrete counting question is substantial on its own.}
\labelled{Ten crossings and its limits.}{For integer dimensions, reduce to coprime $u,v$ and solve $(u-1)(v-1)=20$. The positive factor pairs give reduced rectangles 2-by-21, 3-by-11, 5-by-6 and their swaps. Only the latter two fit the 12-by-12 workspace; a double of 5-by-6 also fits. The independent finite checker verifies all 144 integer rectangles within that workspace. Other similar copies have the same count when drawn as continuous ideal paths; if children scale to noninteger side lengths, use a matching smaller unit rather than forcing unit counter steps through a wall.}
\labelled{Further work.}{Ask whether a path ever follows a previous segment backwards before its terminal corner. In this model it does not: an overlapping same-slope segment would force an earlier repeated reflected state; the coprime reduced path contains neither. Changed slopes or interior starts are different questions and need fresh checks. Save circles and drawings so the next visit can compare counts without retracing every table.}

\end{document}
